% c16 原书 p.195–210（PDF p.202–217）
% 起始页 p.202 顶部为 §2.4 习题 (1)–(4) 的续文（属 c15），按接缝规则跳过，自本节标题起译。
\subsection{固有态射的上同调研究}
\label{sec:2.5}

设 $A$ 是诺特环，$X$ 是 $\Spec A$ 上的固有概形，$\cF$ 是凝聚 $\cO_X$-模。本节我们证明：对每个 $n$，$\HH^n(X,\cF)$ 是有限生成 $A$-模。此外，对 $A$ 的任意理想 $I$，我们证明
\[
(\HH^n(X,\cF))^{\wedge}\cong \dlim_k \HH^n(X,\cF/I^k\cF),
\]
其中 $(\HH^n(X,\cF))^{\wedge}=\dlim_k \HH^n(X,\cF)/I^k\HH^n(X,\cF)$ 是 $\HH^n(X,\cF)$ 的 $I$-进完全化。随后我们用这些结果得到固有态射的 \emph{Stein 分解}\idx{steinfenjie}{Stein 分解}、连通性定理\idx{liantongxingdingli}{连通性定理}与 \emph{Zariski 主定理}\idx{zariskizhudingli}{Zariski 主定理}。

\begin{theorem}\label{it:2.5.1}
设 $f\colon X\to Y$ 是诺特概形的固有态射。则对每个凝聚 $\cO_X$-模 $\cF$ 与每个 $n$，$R^nf_*\cF$ 是凝聚 $\cO_Y$-模。
\end{theorem}
\begin{proof}
由 \ref{it:2.4.9} 已知 $R^nf_*\cF$ 是拟凝聚的，故只需证明它是有限型的。令 $\mathcal{S}$ 为满足下述条件的 $X$ 的闭子集 $Z$ 的族：存在凝聚 $\cO_X$-模 $\cF$，使 $\Supp\cF\subset Z$，且对某个 $n$，$R^nf_*\cF$ 不凝聚。若 $\mathcal{S}=\emptyset$，则证毕。设 $\mathcal{S}\neq\emptyset$。由于 $X$ 诺特，$\mathcal{S}$ 有极小元。取 $Z$ 为 $\mathcal{S}$ 的极小元，取凝聚 $\cO_X$-模 $\cF$ 使 $\Supp\cF\subset Z$ 且对某个 $n$，$R^nf_*\cF$ 不凝聚。

在 $Z$ 上取诱导既约闭子概形结构，令 $i\colon Z\to X$ 为闭浸入。我们首先说明，可以取 $\cF$ 形如 $\cF=i_*\cG$，其中 $\cG$ 是某个凝聚 $\cO_Z$-模。设 $\mathcal{I}$ 是闭子概形 $Z$ 的理想层。我们断言：存在 $m\ge 0$，使 $R^nf_*(\mathcal{I}^m\cF/\mathcal{I}^{m+1}\cF)$ 对某个 $n$ 不凝聚。于是取 $\cG=i^*(\mathcal{I}^m\cF/\mathcal{I}^{m+1}\cF)$。它是凝聚 $\cO_Z$-模，且 $i_*\cG\cong\mathcal{I}^m\cF/\mathcal{I}^{m+1}\cF$。因此 $R^nf_*(i_*\cG)\cong R^nf_*(\mathcal{I}^m\cF/\mathcal{I}^{m+1}\cF)$ 对某个 $n$ 不凝聚。下面用反证法证明断言。设对一切 $m\ge 0$ 与一切 $n$，$R^nf_*(\mathcal{I}^m\cF/\mathcal{I}^{m+1}\cF)$ 都凝聚。取 $m=0$ 可见，对一切 $n$，$R^nf_*(\cF/\mathcal{I}\cF)$ 凝聚。设已证对一切 $n$，$R^nf_*(\cF/\mathcal{I}^T\cF)$ 凝聚。有正合序列
\[
R^nf_*(\mathcal{I}^T\cF/\mathcal{I}^{T+1}\cF)\to R^nf_*(\cF/\mathcal{I}^{T+1}\cF)\to R^nf_*(\cF/\mathcal{I}^T\cF).
\]
这说明对一切 $n$，$R^nf_*(\cF/\mathcal{I}^{T+1}\cF)$ 凝聚。故对一切 $n$ 与一切 $m\ge 0$，$R^nf_*(\cF/\mathcal{I}^m\cF)$ 都凝聚。由于 $\Supp\cF\subset Z$，由 \ref{it:1.4.14} 有 $\mathcal{I}^k\cF=0$ 对某个自然数 $k$ 成立。于是对一切 $n$，$R^nf_*\cF=R^nf_*(\cF/\mathcal{I}^k\cF)$ 凝聚。这与 $\cF$ 的取法矛盾。故断言成立。

若 $Z$ 不是不可约的，则 $Z=Z_1\cup Z_2$，其中 $Z_1$、$Z_2$ 是 $Z$ 的真闭子集。在 $Z_1$、$Z_2$ 上取诱导既约闭子概形结构，令 $i_1\colon Z_1\to Z$、$i_2\colon Z_2\to Z$ 为闭浸入，并令
\[
\phi\colon i_*\cG\to i_*i_{1*}i_1^*\cG\oplus i_*i_{2*}i_2^*\cG
\]
为由典范态射 $\cG\to i_{1*}i_1^*\cG$ 与 $\cG\to i_{2*}i_2^*\cG$ 定义的态射。注意 $\phi$ 在 $Z_1\cap Z_2$ 之外是同构。故 $\ker\phi$ 与 $\coker\phi$ 的支集落在 $Z_1\cap Z_2$ 中。由于 $Z$ 是 $\mathcal{S}$ 的极小元，$Z_1\cap Z_2$、$Z_1$ 与 $Z_2$ 都不在 $\mathcal{S}$ 中。故对一切 $n$，$\cO_Y$-模 $R^nf_*(\ker\phi)$、$R^nf_*(\coker\phi)$、$R^nf_*(i_*i_{1*}i_1^*\cG)$ 与 $R^nf_*(i_*i_{2*}i_2^*\cG)$ 都凝聚。有正合序列
\[
\cdots\to R^{n-1}f_*(\coker\phi)\to R^nf_*(\operatorname{im}\phi)\to R^nf_*(i_*i_{1*}i_1^*\cG\oplus i_*i_{2*}i_2^*\cG)\to\cdots.
\]
这说明对一切 $n$，$R^nf_*(\operatorname{im}\phi)$ 凝聚。又有正合序列
\[
\cdots\to R^nf_*(\ker\phi)\to R^nf_*(i_*\cG)\to R^nf_*(\operatorname{im}\phi)\to\cdots.
\]
这说明对一切 $n$，$R^nf_*(i_*\cG)=R^nf_*\cF$ 凝聚。这与 $\cF$ 的取法矛盾。

设 $Z$ 不可约。令 $z$ 为 $Z$ 的\emph{一般点}。由 Chow 引理 \ref{it:1.4.18}，存在射影态射 $g\colon Z'\to Z$，使 $fig$ 也是射影的，且存在 $Z$ 的稠密开子集 $U$，使 $g$ 诱导出 $g^{-1}(U)$ 与 $U$ 的同构。注意 $z$ 落在 $U$ 中。令 $\cO_{Z'}(1)$ 为关于 $Z$ 极丰沛的可逆 $\cO_{Z'}$-模。由 \ref{it:2.4.13}(iii)，可取充分大的 $k$，使对一切 $q\ge 1$ 有 $R^qg_*\cO_{Z'}(k)=0$。于是双正则谱序列
\[
R^p(fi)_*R^qg_*\cO_{Z'}(k)\Rightarrow R^{p+q}(fig)_*\cO_{Z'}(k)
\]
退化，从而对一切 $n$ 有
\[
R^n(fi)_*(g_*\cO_{Z'}(k))\cong R^n(fig)_*\cO_{Z'}(k).
\]
利用 $i_*$ 是正合函子且把 $Z$ 上的内射层映为 $X$ 上的内射层这一事实，可以验证
\[
R^nf_*(i_*g_*\cO_{Z'}(k))\cong R^n(fi)_*(g_*\cO_{Z'}(k)).
\]
由 \ref{it:2.4.13}(i)，对一切 $n$，$g_*\cO_{Z'}(k)$ 与 $R^n(fig)_*\cO_{Z'}(k)$ 都凝聚。故对一切 $n$，$R^nf_*(i_*g_*\cO_{Z'}(k))$ 都凝聚。由于 $g$ 诱导出 $g^{-1}(U)$ 与 $U$ 的同构，我们有 $(g_*\cO_{Z'}(k))_z\cong\cO_{Z,z}$。注意 $\cG_z$ 是域 $\cO_{Z,z}$ 上的有限维向量空间。故存在自然数 $m$，使有 $\cO_{Z,z}$-向量空间的同构
\[
\phi_z\colon (g_*\cO_{Z'}(k))_z^m\cong\cG_z.
\]
由 \ref{it:1.4.1}(ii)，存在 $z$ 在 $X$ 中的开邻域 $V$，使有同构
\[
\phi_V\colon (i_*g_*\cO_{Z'}(k))^m|_V\cong i_*\cG|_V,
\]
它在茎上诱导出同构 $\phi_z$。令 $\mathcal{H}_V$ 为 $\phi_V$ 的图像，即态射
\[
(i_*g_*\cO_{Z'}(n))^m|_V\to\bigl((i_*g_*\cO_{Z'}(k))^m\oplus i_*\cG\bigr)|_V,\qquad s\mapsto(s,\phi_V(s))
\]
的像。

注意 $\mathcal{H}_V$ 是凝聚 $\cO_V$-模。由下面的 \ref{it:2.5.2}，存在凝聚 $\cO_X$-子模 $\mathcal{H}\subset (i_*g_*\cO_{Z'}(k))^m\oplus i_*\cG$，使 $\mathcal{H}|_V\cong\mathcal{H}_V$ 且 $\mathcal{H}|_{X-Z}=0$。考虑投影 $p_1\colon\mathcal{H}\to(i_*g_*\cO_{Z'}(k))^m$ 与 $p_2\colon\mathcal{H}\to i_*\cG$。它们的核与余核的支集落在 $Z-V$ 中。由于 $Z$ 是 $\mathcal{S}$ 的极小元，对一切 $n$，$R^nf_*(\ker p_1)$、$R^nf_*(\coker p_1)$、$R^nf_*(\ker p_2)$ 与 $R^nf_*(\coker p_2)$ 都凝聚。结合 $R^nf_*((i_*g_*\cO_{Z'}(k))^m)$ 对一切 $n$ 凝聚这一事实，我们得到：对一切 $n$，$R^nf_*(i_*\cG)$ 凝聚。故对一切 $n$，$R^nf_*\cF=R^nf_*(i_*\cG)$ 凝聚。这与 $\cF$ 的取法矛盾。

无论哪种情形，假设 $\mathcal{S}\neq\emptyset$ 都导致矛盾。故 $\mathcal{S}=\emptyset$。
\end{proof}

\begin{tnote}
上式中第一项的 $\cO_{Z'}(n)$ 与原书其余各处的 $\cO_{Z'}(k)$ 不一致：该态射的出发项与到达项的第一分量应当是同一个模 $(i_*g_*\cO_{Z'}(k))^m|_V$，否则「图像」无从谈起，故此处 $n$ 应为 $k$ 之误植。译文照原书保留 $n$，特此说明。[依据：原书 p.197 定理 2.5.1 证明中 $\mathcal{H}_V$ 的定义]
\end{tnote}
\noindent\textbf{注}
我们由 \ref{it:2.4.13} 给出 \ref{it:2.5.1} 的另一个推导。沿用上面的记号。令
\[
0\to\cG\to\cC^\bullet(\cG)\quad(\text{resp. } 0\to g^*\cG\to\cC^\bullet(g^*\cG))
\]
为 $\cG$（resp.\ $g^*\cG$）的 Godement 分解，令
\[
\phi\colon\cC^\bullet(\cG)\to g_*(\cC^\bullet(g^*\cG))
\]
为由 $\cG\to g_*g^*\cG$ 诱导的复形态射，并令 $\cC^\bullet$ 为 $\phi$ 的映射锥。我们有长正合序列
\[
\cdots\to\cH^q(\cC^\bullet(\cG))\to\cH^q(g_*(\cC^\bullet(g^*\cG)))\to\cH^q(\cC^\bullet)\to\cH^{q+1}(\cC^\bullet(\cG))\to\cdots.
\]
于是有同构
\[
R^qg_*g^*\cG\cong\cH^q(\cC^\bullet)
\]
对 $q\neq 0,-1$ 成立，以及正合序列
\[
0\to\cH^{-1}(\cC^\bullet)\to\cG\to g_*g^*\cG\to\cH^0(\cC^\bullet)\to 0.
\]
由 \ref{it:2.4.13}，对每个 $q$，$R^qg_*g^*\cG$ 凝聚。故对每个 $q$，$\cH^q(\cC^\bullet)$ 凝聚。由于 $g$ 在 $U$ 上是同构，我们有
\[
\cG|_U\cong(g_*g^*\cG)|_U,\qquad R^qg_*g^*\cG=0\quad(q\ge 1).
\]
故对每个 $q$，$\cH^q(\cC^\bullet)|_U=0$。由 $Z$ 的极小性，对每对 $p,q$，$R^p(fi)_*\cH^q(\cC^\bullet)$ 凝聚。注意 $\cC^\bullet$ 的各项是松弛层。令 $\cI^\bullet$ 为 $\cC^\bullet$ 的 Cartan-Eilenberg 分解。通过分析双复形 $(fi)_*\cI^\bullet$ 关联的两个谱序列，我们得到谱序列
\[
E_2^{pq}=R^p(fi)_*\cH^q(\cC^\bullet)\Rightarrow \HH^{p+q}(fi)_*\cC^\bullet.
\]
由此可知，对每个 $n$，$\cH^n(fi)_*\cC^\bullet$ 凝聚。注意 $(fi)_*\cC^\bullet$ 是态射
\[
(fi)_*\phi\colon(fi)_*\cC^\bullet(\cG)\to(fig)_*(\cC^\bullet(g^*\cG))
\]
的映射锥。故有长正合序列
\[
\cdots\to\cH^{n-1}(fi)_*\cC^\bullet\to R^n(fi)_*\cG\to R^n(fig)_*g^*\cG\to\cH^n(fi)_*\cC^\bullet\to\cdots.
\]
由 \ref{it:2.4.13}，对每个 $n$，$R^n(fig)_*g^*\cG$ 拟凝聚。于是我们断定：对每个 $n$，$R^n(fi)_*\cG$ 拟凝聚。矛盾。

\begin{lemma}\label{it:2.5.2}
设 $X$ 是诺特概形，$U$ 是 $X$ 的开子概形，$\cF$ 是拟凝聚 $\cO_X$-模，$\cG_U$ 是 $\cF|_U$ 的凝聚 $\cO_U$-子模。则存在 $\cF$ 的凝聚 $\cO_X$-子模 $\cG$，使 $\cG|_U=\cG_U$。
\end{lemma}
\begin{proof}
先考虑 $X=\Spec A$ 是仿射的情形。令 $j\colon U\to X$ 为开浸入。由于 $\cG_U$ 是 $j^*\cF=\cF|_U$ 的子模，$j_*\cG_U$ 是 $j_*j^*\cF$ 的子模。令 $\cG'$ 为 $j_*\cG_U$ 在典范态射 $\cF\to j_*j^*\cF$ 下的原像。由于 $\cF$、$j_*j^*\cF$ 与 $j_*\cG_U$ 由 \ref{it:1.4.9} 都拟凝聚，$\cG'$ 拟凝聚。注意 $j^*\cG'$ 是 $j^*j_*\cG_U=\cG_U$ 在恒等态射
\[
j^*\cF\to j^*(j_*j^*\cF)=j^*\cF
\]
下有原像。故 $j^*\cG'=\cG|_U$。
\begin{tnote}
原书此式印作 $j^*\cG'=\cG|_U$，但 $\cG$ 在本书此前的记号中并未定义，疑本应写 $\cG_U$（即 $\cG$ 在 $U$ 上的限制）；译文照录原书印法。[依据：原书 p.199 第 2 段]
\end{tnote}
设 $\cG'\cong M^{\sim}$，其中 $M$ 是 $A$-模。我们有 $M=\dlim_n M_n$，其中有向极限取遍 $M$ 的有限生成 $A$-子模 $M_n$ 构成的有向系统。于是 $\cG'\cong\dlim_n\cG_n$，其中每个 $\cG_n=M_n^{\sim}$ 都是 $\cF$ 的凝聚 $\cO_X$-子模。注意
\[
\cG_U\cong j^*\cG'\cong\dlim_n j^*\cG_n.
\]
我们断言：对某个 $n$ 有 $\cG_U\cong j^*\cG_n$。这将证明 $X$ 仿射情形的引理。用有限多个仿射开子概形 $U_i=\Spec A_i$（$i\in I$）覆盖 $U$。由于 $\cG_U$ 是凝聚 $\cO_U$-模，有 $\cG_U|_{U_i}\cong N_i^{\sim}$，其中 $N_i$ 是某个有限生成 $A_i$-模。我们有
\[
\cG_U|_{U_i}\cong\dlim_n\cG_n|_{U_i}=\dlim_n (A_i\tensor_A M_n)^{\sim}.
\]
故 $N_i\cong\dlim_n A_i\tensor_A M_n$。由于每个 $N_i$ 都是有限生成 $A_i$-模，存在 $M_n$，使对每个 $i\in I$ 有 $N_i\cong A_i\tensor_A M_n$。于是得到 $\cG_U\cong j^*\cG_n$。

一般情形，用有限多个仿射开子概形 $U_1,\dots,U_n$ 覆盖 $X$。令 $V_i=U\cup U_1\cup\cdots\cup U_i$（$1\le i\le n$）。我们对 $i$ 归纳构造 $\cF|_{V_i}$ 的凝聚 $\cO_{V_i}$-子模 $\cG_{V_i}$，使 $\cG_{V_i}|_U\cong\cG_U$。于是 $\cG=\cG_{V_n}$ 具有所要的性质。由上面已处理的仿射情形，存在 $\cF|_{U_1}$ 的凝聚 $\cO_{U_1}$-子模 $\mathcal{H}_{U_1}$，使 $\mathcal{H}_{U_1}|_{U\cap U_1}=\cG_U|_{U\cap U_1}$。我们可以把 $\mathcal{H}_{U_1}$ 与 $\cG_U$ 粘成 $V_1=U\cup U_1$ 上的层 $\cG_{V_1}$。设已构造出 $\cG_{V_i}$。由上面已处理的仿射情形，可找到 $\cF|_{U_{i+1}}$ 的凝聚 $\cO_{U_{i+1}}$-子模 $\mathcal{H}_{U_{i+1}}$，使 $\mathcal{H}_{U_{i+1}}|_{V_i\cap U_{i+1}}=\cG_{V_i}|_{V_i\cap U_{i+1}}$。我们可以把 $\mathcal{H}_{U_{i+1}}$ 与 $\cG_{V_i}$ 粘成 $V_i\cup U_{i+1}=V_{i+1}$ 上的层 $\cG_{V_{i+1}}$。
\end{proof}

\begin{corollary}[有限性定理]\label{it:2.5.3}
设 $A$ 是诺特环，$X$ 是 $\Spec A$ 上的固有概形，$\cF$ 是凝聚 $\cO_X$-模。则对一切 $n$，$\HH^n(X,\cF)$ 都是有限生成 $A$-模。
\end{corollary}

\begin{proof}
由 \ref{it:2.5.1} 与 \ref{it:2.4.9} 即得。
\end{proof}

\begin{corollary}\label{it:2.5.4}
(i) 设 $X$ 是诺特概形，$\cF$ 与 $\cG$ 是凝聚 $\cO_X$-模。则对一切 $n$，$\Ext^n_{\cO_X}(\cF,\cG)$ 是凝聚 $\cO_X$-模。

(ii) 设 $A$ 是诺特环，$X$ 是 $\Spec A$ 上的固有概形，$\cF$ 与 $\cG$ 是凝聚 $\cO_X$-模。则对一切 $n$，$\Ext^n_{\cO_X}(\cF,\cG)$ 都是有限生成 $A$-模。
\end{corollary}
\begin{proof}
(i) 这是局部问题。可设 $X=\Spec A$，其中 $A$ 是诺特环。则 $\cF\cong M^{\sim}$，其中 $M$ 是有限生成 $A$-模。由于 $M$ 有由有限秩自由 $A$-模构成的分解，我们可以找到 $\cF$ 的由有限秩自由 $\cO_X$-模构成的分解
\[
\cdots\to\cE_1\to\cE_0\to\cF\to 0.
\]
由 \ref{it:2.2.16}，我们有
\[
\Ext^n_{\cO_X}(\cF,\cG)=\HH^n(\mathcal{H}om_{\cO_X}(\cE,\cG)).
\]
注意每个 $\mathcal{H}om_{\cO_X}(\cE_n,\cG)$ 都凝聚。故每个 $\Ext^n_{\cO_X}(\cF,\cG)$ 也凝聚。

(ii) 我们有双正则谱序列
\[
\HH^p(X,\Ext^q_{\cO_X}(\cF,\cG))\Rightarrow\Ext^{p+q}_{\cO_X}(\cF,\cG).
\]
由 \ref{it:2.5.3}，每个 $\HH^p(X,\Ext^q_{\cO_X}(\cF,\cG))$ 都是有限生成 $A$-模。这说明每个 $\Ext^n_{\cO_X}(\cF,\cG)$ 都是有限生成 $A$-模。
\end{proof}

\begin{lemma}\label{it:2.5.5}
设 $A$ 是诺特环，$I$ 是 $A$ 的理想，$S$ 是分次环 $\bigoplus_{k=0}^\infty I^k$，$f\colon X\to\Spec A$ 是固有态射，$\cF$ 是凝聚 $\cO_X$-模。则 $S$ 是诺特的，且对每个 $n$，$\bigoplus_{k=0}^\infty \HH^n(X,I^k\cF)$ 是有限生成分次 $S$-模，其中对一切 $k$ 有 $I^k\cF=(I^k)^{\sim}\cF$。
\end{lemma}

\begin{proof}
设 $a_1,\dots,a_m$ 是 $I$ 的有限生成元组。$A$-代数同态
\[
A[x_1,\dots,x_m]\to S,\qquad x_i\mapsto a_i\in S_1
\]
是满射。由于 $A[x_1,\dots,x_m]$ 诺特，$S$ 也诺特。对每个 $a\in S_k=I^k$，乘以 $a$ 诱导出态射 $I^l\cF\to I^{k+l}\cF$，从而诱导出同态 $\HH^n(X,I^l\cF)\to\HH^n(X,I^{k+l}\cF)$。这就赋予 $\bigoplus_{k=0}^\infty \HH^n(X,I^k\cF)$ 以分次 $S$-模结构。

用下面的笛卡尔图确定记号：
\[
\begin{tikzcd}
X'=X\fib_{\Spec A}\Spec S \arrow[r,"g'"] \arrow[d,"f'"'] & X \arrow[d,"f"] \\
\Spec S \arrow[r,"g"'] & \Spec A
\end{tikzcd}
\]
其中 $g$ 对应同态
\[
A\hookrightarrow S=A\oplus I\oplus I^2\oplus\cdots.
\]
注意 $g$ 与 $g'$ 都是仿射态射，且
\[
g'_*\cO_{X'}\cong f^*(S^{\sim})\cong\bigoplus_{k=0}^\infty f^*((I^k)^{\sim}).
\]
$X$ 上的层 $\bigoplus_{k=0}^\infty I^k\cF$ 是 $\bigoplus_{k=0}^\infty f^*((I^k)^{\sim})$-模，作为 $\cO_X$-模是拟凝聚的。由 \ref{it:1.4.12}，存在拟凝聚 $\cO_{X'}$-模 $\cF'$，使 $g'_*\cF'\cong\bigoplus_{k=0}^\infty I^k\cF$。注意 $\cF'$ 实际上是凝聚的。由于 $g'$ 仿射，双正则谱序列
\[
\HH^p(X,R^qg'_*\cF')\Rightarrow\HH^{p+q}(X',\cF')
\]
由 \ref{it:2.4.8} 退化。于是有
\[
\HH^n(X',\cF')\cong\HH^n(X,g'_*\cF')\cong\HH^n\Bigl(X,\bigoplus_{k=0}^\infty I^k\cF\Bigr)\cong\bigoplus_{k=0}^\infty\HH^n(X,I^k\cF).
\]
这是分次 $S$-模的同构，只要用仿射开覆盖的 Čech 上同调来计算这些上同调群即可看出。由 \ref{it:2.5.3}，$\HH^n(X',\cF')$ 是有限生成 $S$-模。故 $\bigoplus_{k=0}^\infty\HH^n(X,I^k\cF)$ 是有限生成分次 $S$-模。
\end{proof}

\begin{theorem}\label{it:2.5.6}
设 $A$ 是诺特环，$I$ 是 $A$ 的理想，$X$ 是 $\Spec A$ 上的固有概形，$\cF$ 是凝聚 $\cO_X$-模。则对每个 $n$ 有
\[
(\HH^n(X,\cF))^{\wedge}\cong\dlim_k \HH^n(X,\cF/I^k\cF),
\]
其中 $(\HH^n(X,\cF))^{\wedge}=\dlim_k \HH^n(X,\cF)/I^k\HH^n(X,\cF)$ 是 $A$-模 $\HH^n(X,\cF)$ 的 $I$-进完全化。
\end{theorem}

\begin{proof}
对每个 $n$ 与每个 $k$，有 $I^k\HH^n(X,\cF/I^k\cF)=0$。故典范同态 $\HH^n(X,\cF)\to\HH^n(X,\cF/I^k\cF)$ 诱导出同态
\[
\HH^n(X,\cF)/I^k\HH^n(X,\cF)\to\HH^n(X,\cF/I^k\cF).
\]
取逆向极限，得到同态
\[
(\HH^n(X,\cF))^{\wedge}\to\dlim_k\HH^n(X,\cF/I^k\cF)
\]
对每个 $n$。下面证明它是同构。

我们有长正合序列
\[
\cdots\to\HH^n(X,I^k\cF)\to\HH^n(X,\cF)\to\HH^n(X,\cF/I^k\cF)\to\HH^{n+1}(X,I^k\cF)\to\HH^{n+1}(X,\cF)\to\cdots.
\]
令
\[
\begin{aligned}
R_k&=\operatorname{im}\bigl(\HH^n(X,I^k\cF)\to\HH^n(X,\cF)\bigr),\\
Q_k&=\ker\bigl(\HH^{n+1}(X,I^k\cF)\to\HH^{n+1}(X,\cF)\bigr)\\
&=\operatorname{im}\bigl(\HH^n(X,\cF/I^k\cF)\to\HH^{n+1}(X,I^k\cF)\bigr).
\end{aligned}
\]
于是有正合序列
\[
0\to\HH^n(X,\cF)/R_k\to\HH^n(X,\cF/I^k\cF)\to Q_k\to 0.
\]
取逆向极限并应用 \ref{it:1.5.1}，得到正合序列
\[
0\to\dlim_k\HH^n(X,\cF)/R_k\to\dlim_k\HH^n(X,\cF/I^k\cF)\to\dlim_k Q_k\to 0.
\]
由下面的 \ref{it:2.5.7}，有 $\dlim_k\HH^n(X,\cF)/R_k\cong(\HH^n(X,\cF))^{\wedge}$。由下面的 \ref{it:2.5.8}，有 $\dlim_k Q_k=0$。于是得到
\[
(\HH^n(X,\cF))^{\wedge}\cong\dlim_k\HH^n(X,\cF/I^k\cF).
\]
\end{proof}

\begin{lemma}\label{it:2.5.7}
对每个 $n$，由 $R_k$ 给出的 $\HH^n(X,\cF)$ 上的递减滤链是 $I$-稳定的，即对一切 $k,l\ge 0$ 有 $I^lR_k\subset R_{k+l}$，且存在自然数 $k_0$，使对一切 $k\ge k_0$ 与 $l\ge 0$ 有 $I^lR_k=R_{k+l}$。此外，我们有
\[
(\HH^n(X,\cF))^{\wedge}\cong\dlim_k\HH^n(X,\cF)/R_k.
\]
\end{lemma}

\begin{proof}
对每个 $a\in I^l$，有交换图
\[
\begin{tikzcd}
I^k\cF \arrow[r,"a"] \arrow[d] & I^{k+l}\cF \arrow[d] \\
\cF \arrow[r,"a"'] & \cF
\end{tikzcd}
\]
其中水平箭头由乘以 $a$ 诱导，竖直箭头是包含。这个图诱导出
\[
\begin{tikzcd}
\HH^n(X,I^k\cF) \arrow[r,"a"] \arrow[d] & \HH^n(X,I^{k+l}\cF) \arrow[d] \\
\HH^n(X,\cF) \arrow[r,"a"'] & \HH^n(X,\cF)
\end{tikzcd}
\]
故乘以 $a$ 把第一个竖直箭头的像映入第二个竖直箭头的像，即 $I^lR_k\subset R_{k+l}$。这使得 $\bigoplus_{k=0}^\infty R_k$ 成为分次 $S=\bigoplus_{k=0}^\infty I^k$-模。它是分次 $S$-模的商，而后者为 $\bigoplus_{k=0}^\infty\HH^n(X,I^k\cF)$。由 \ref{it:2.5.5}，$S$ 诺特且 $\bigoplus_{k=0}^\infty\HH^n(X,I^k\cF)$ 有限生成。故 $\bigoplus_{k=0}^\infty R_k$ 是有限生成分次 $S$-模。设 $x_i$（$i=1,\dots,m$）是生成 $\bigoplus_{k=0}^\infty R_k$ 的有限个齐次元素，令 $k_0=\max(\deg x_1,\dots,\deg x_m)$。可以验证，对每个 $l\ge 0$ 有 $R_{k_0+l}=I^lR_{k_0}$。当 $k\ge k_0$ 时，我们有
\[
R_k=I^{k-k_0}R_{k_0}\subset I^{k-k_0}\HH^n(X,\cF).
\]
此外，对每个 $k\ge 0$ 有
\[
I^k\HH^n(X,\cF)=I^kR_0\subset R_k.
\]
故
\[
(\HH^n(X,\cF))^{\wedge}=\dlim_k\HH^n(X,\cF)/I^k\HH^n(X,\cF)\cong\dlim_k\HH^n(X,\cF)/R_k.
\]
\end{proof}
\begin{lemma}\label{it:2.5.8}
存在整数 $k_0$，使典范同态 $Q_{k+k_0}\to Q_k$ 对每个 $k\ge k_0$ 都为零。
\end{lemma}

\begin{proof}
对每个 $a\in I^l$，有交换图
\[
\begin{tikzcd}
\HH^{n+1}(X,I^k\cF) \arrow[r,"a"] \arrow[d] & \HH^{n+1}(X,I^{k+l}\cF) \arrow[d] \\
\HH^{n+1}(X,\cF) \arrow[r,"a"'] & \HH^{n+1}(X,\cF)
\end{tikzcd}
\]
其中水平箭头由乘以 $a$ 诱导。故乘以 $a$ 把第一个竖直箭头的核映入第二个竖直箭头的核，即 $I^lQ_k\subset Q_{k+l}$。这使得 $\bigoplus_{k=0}^\infty Q_k$ 成为分次 $S=\bigoplus_{k=0}^\infty I^k$-模，它是 $\bigoplus_{k=0}^\infty\HH^{n+1}(X,I^k\cF)$ 的分次子模。由 \ref{it:2.5.5}，$S$ 诺特且 $\bigoplus_{k=0}^\infty\HH^{n+1}(X,I^k\cF)$ 有限生成。故 $\bigoplus_{k=0}^\infty Q_k$ 是有限生成分次 $S$-模。设 $x_1,\dots,x_m$ 是生成 $\bigoplus_{k=0}^\infty Q_k$ 的有限个齐次元素，令 $k_0=\max(\deg x_1,\dots,\deg x_m)$。可以验证，对一切 $k\ge k_0$ 与 $l\ge 0$ 有 $Q_{k+l}=S_lQ_k$。另一方面，由于 $A$-模 $\HH^n(X,\cF/I^k\cF)$ 被 $I^k$ 零化，且 $Q_k=\operatorname{im}(\HH^n(X,\cF/I^k\cF)\to\HH^{n+1}(X,I^k\cF))$，$A$-模 $Q_k$ 也被 $I^k$ 零化，即
\[
\alpha_k(S_k)Q_k=0
\]
对每个 $k\ge 0$ 成立，其中 $\alpha_k\colon S_k\to S_0$ 是包含 $I^k\hookrightarrow A$。这说明
\[
\alpha_{k_0}(S_{k_0})x_i=0\qquad(i=1,\dots,m).
\]
故
\[
\alpha_{k_0}(S_{k_0})Q=0.
\]
对每个 $a\in S_{k_0}$，有交换图
\[
\begin{tikzcd}
Q_k \arrow[r] \arrow[d] & Q_{k+k_0} \arrow[r] \arrow[d] & Q_k \arrow[d] \\
\HH^{n+1}(X,I^k\cF) \arrow[r,"a"] & \HH^{n+1}(X,I^{k+k_0}\cF) \arrow[r,"\alpha_k(a)"'] & \HH^{n+1}(X,I^k\cF)
\end{tikzcd}
\]
对每个 $k\ge k_0$，由于 $Q_{k+k_0}=S_{k_0}Q_k$，$Q_{k+k_0}\to Q_k$ 的像就是 $\alpha_{k_0}(S_{k_0})Q_k$。故 $Q_{k+k_0}\to Q_k$ 对每个 $k\ge k_0$ 都为零。
\end{proof}

\begin{corollary}\label{it:2.5.9}
设 $f\colon X\to Y$ 是诺特概形的固有态射，$\cF$ 是凝聚 $\cO_X$-模，$\mathcal{I}$ 是 $\cO_Y$ 的凝聚理想层。则对每个 $n$，有典范同构
\[
\dlim_k R^nf_*\cF/\mathcal{I}^kR^nf_*\cF\cong\dlim_k R^nf_*(\cF/\mathcal{I}^k\cF).
\]
\end{corollary}
\begin{proof}
注意，典范态射 $R^nf_*\cF\to R^nf_*(\cF/\mathcal{I}^k\cF)$ 对每个 $n$ 与每个 $k$ 诱导出态射
\[
R^nf_*\cF/\mathcal{I}^kR^nf_*\cF\to R^nf_*(\cF/\mathcal{I}^k\cF).
\]
取逆向极限，对每个 $n$ 得到态射
\[
\dlim_k R^nf_*\cF/\mathcal{I}^kR^nf_*\cF\to\dlim_k R^nf_*(\cF/\mathcal{I}^k\cF).
\]
我们来证明它是同构。只需证明
\[
\dlim_k(R^nf_*\cF/\mathcal{I}^kR^nf_*\cF)(V)\to\dlim_k R^nf_*(\cF/\mathcal{I}^k\cF)(V)
\]
对每个仿射开子概形 $V=\Spec A\subset Y$ 都是同构。设 $I$ 是 $A$ 中对应于 $\mathcal{I}$ 的理想。则由 \ref{it:2.4.9}，我们有
\[
(R^nf_*\cF/\mathcal{I}^kR^nf_*\cF)(V)\cong\HH^n(f^{-1}(V),\cF)/I^k\HH^n(f^{-1}(V),\cF),
\]
\[
R^nf_*(\cF/\mathcal{I}^k\cF)(V)\cong\HH^n(f^{-1}(V),\cF/\mathcal{I}^k\cF).
\]
于是我们的断言由 \ref{it:2.5.6} 推出。
\end{proof}

\begin{theorem}[连通性定理]\label{it:2.5.10}
设 $f\colon X\to Y$ 是诺特概形的固有态射，且 $f^\sharp\colon\cO_Y\to f_*\cO_X$ 是同构。则对每个 $y\in Y$，$f^{-1}(y)$ 非空且连通。
\end{theorem}

\begin{proof}
先证 $f$ 是满射，从而 $f^{-1}(y)$ 非空。由于 $f$ 固有，$f(X)$ 是闭集。若 $f(X)\neq Y$，取 $y\in Y-f(X)$，则 $(f_*\cO_X)_y=0$。但 $(f_*\cO_X)_y\cong\cO_{Y,y}\neq 0$，矛盾。故 $f(X)=Y$。

存在典范的平坦态射 $\Spec\cO_{Y,y}\to Y$。事实上，设 $V\cong\Spec A$ 是 $y$ 在 $Y$ 中的仿射开邻域，$\mathfrak{p}$ 是 $A$ 中对应于 $y$ 的素理想。则它是复合
\[
\Spec\cO_{Y,y}\cong\Spec A_{\mathfrak{p}}\to\Spec A\cong V\hookrightarrow Y.
\]
该复合与仿射开邻域 $V$ 的选取无关。用下面的笛卡尔图固定记号：
\[
\begin{tikzcd}
X_k=X\times_Y\Spec\cO_{Y,y}/\fm_y^k \arrow[r,"i_k"] \arrow[d,"f_k"] & X'=X\times_Y\Spec\cO_{Y,y} \arrow[r,"i"] \arrow[d,"f'"] & X \arrow[d,"f"] \\
\Spec\cO_{Y,y}/\fm_y^k \arrow[r] & \Spec\cO_{Y,y} \arrow[r,"j"] & Y
\end{tikzcd}
\]
其中 $\fm_y$ 是 $\cO_{Y,y}$ 的极大理想。由 \ref{it:2.4.10}，我们有
\[
f'_*\cO_{X'}\cong f'_*i^*\cO_X\cong j^*f_*\cO_X\cong j^*\cO_Y\cong\cO_{\Spec\cO_{Y,y}}.
\]
故有
\[
\HH^0(X',\cO_{X'})\cong\Gamma(\Spec\cO_{Y,y},f'_*\cO_{X'})\cong\Gamma(\Spec\cO_{Y,y},\cO_{\Spec\cO_{Y,y}})\cong\cO_{Y,y}.
\]

取 $\fm_y$-进完全化，得到
\[
(\HH^0(X',\cO_{X'}))^{\wedge}\cong\hat{\cO}_{Y,y}.
\]
结合 \ref{it:2.5.6}，得到
\[
\dlim_k\HH^0(X',\cO_{X'}/\fm_y^k\cO_{X'})\cong\hat{\cO}_{Y,y}.
\]
我们有
\[
\cO_{X'}/\fm_y^k\cO_{X'}=i_{k*}\cO_{X_k}.
\]
故有
\[
\HH^0(X',\cO_{X'}/\fm_y^k\cO_{X'})\cong\HH^0(X_k,\cO_{X_k}).
\]
因此
\[
\dlim_k\HH^0(X_k,\cO_{X_k})\cong\hat{\cO}_{Y,y}.
\]
特别地，$\dlim_k\HH^0(X_k,\cO_{X_k})$ 是局部环。注意对每个 $k$，$X_k$ 的底拓扑空间与 $f^{-1}(y)$ 相同。若 $f^{-1}(y)$ 不连通，则 $f^{-1}(y)$ 是两个不相交开子集 $U_1$ 与 $U_2$ 的并。我们有
\[
\dlim_k\HH^0(X_k,\cO_{X_k})=\dlim_k\HH^0(U_1,\cO_{X_k})\oplus\dlim_k\HH^0(U_2,\cO_{X_k}).
\]
$\dlim_k\HH^0(U_1,\cO_{X_k})\oplus\dlim_k\HH^0(U_2,\cO_{X_k})$ 中的元素 $(1,0)$ 与 $(0,1)$ 都不是单位元，但它们的和是单位元。故 $\dlim_k\HH^0(X_k,\cO_{X_k})$ 不可能是局部环，矛盾。故 $f^{-1}(y)$ 连通。
\end{proof}

设 $f\colon X\to Y$ 是诺特概形之间的固有态射。由 \ref{it:2.5.1}，$f_*\cO_X$ 是凝聚 $\cO_Y$-模，它是 $\cO_Y$-代数层。由 \ref{it:1.4.12}，存在仿射态射 $g\colon Y'\to Y$，使 $g_*\cO_{Y'}\cong f_*\cO_X$。由 \ref{it:1.4.11}，同构 $g_*\cO_{Y'}\cong f_*\cO_X$ 诱导出态射 $f'\colon X\to Y'$，使 $f$ 可分解为
\[
X\xrightarrow{f'}Y'\xrightarrow{g}Y.
\]
我们称它为 $f$ 的\emph{Stein 分解}。注意 $g$ 是有限态射。由 \ref{it:1.4.10}(ii)，$f'^\sharp\colon\cO_{Y'}\to f'_*\cO_X$ 是同构。故由 \ref{it:2.5.10}，$f'$ 的所有纤维都连通。

称态射 $f\colon X\to Y$ 是\emph{拟有限}\idx{niyouxian}{拟有限}的，若它是有限型的，且对每个 $y\in Y$，$f^{-1}(y)$ 作为集合是有限的。把下面的 \ref{it:2.5.11} 应用于投影
\[
f^{-1}(y)=X\times_Y\Spec k(y)\to\Spec k(y),
\]
其中 $k(y)$ 是局部环 $\cO_{Y,y}$ 的剩余域，我们看到：有限型态射 $f\colon X\to Y$ 拟有限当且仅当对每个 $y\in Y$，$f^{-1}(y)$ 离散。
\begin{lemma}\label{it:2.5.11}
设 $k$ 是域，$X\to\Spec k$ 是有限型态射。则下列条件等价：
\begin{enumerate}[label=(\roman*)]
\item $X$ 的底拓扑空间作为集合是有限的；
\item $X$ 的底拓扑空间是离散的；
\item 存在\emph{阿廷}\idx{ating}{阿廷} $k$-代数 $A$，使 $X\cong\Spec A$。
\end{enumerate}
\end{lemma}

\begin{proof}
(i)$\Rightarrow$(ii) 对每个 $x\in X$，取仿射开邻域 $U=\Spec A$，使 $A$ 是有限生成 $k$-代数。可以证明 $A$ 是 \emph{Jacobson 环}，即 $A$ 的任何素理想都是某些极大理想的交（参见 [Atiyah-Macdonald] 习题 5.24）。由于 $X$ 的底空间是有限集，$A$ 只有有限多个极大理想，设为 $\fm_1,\dots,\fm_n$。若 $\mathfrak{p}=\bigcap_j\fm_j$ 是素理想，则 $\mathfrak{p}\supset\prod_j\fm_j$，从而对某个 $i_j$ 有 $\mathfrak{p}\supset\fm_{i_j}$。由 $\fm_{i_j}$ 的极大性，必有 $\mathfrak{p}=\fm_{i_j}$。故 $A$ 的所有素理想都是极大的。因此有限集 $U$ 中所有点都是闭点。这说明 $U$ 是离散的，从而 $X$ 是离散的。

(ii)$\Rightarrow$(iii) 由于 $X$ 在 $\Spec k$ 上是有限型的，它拟紧。由于 $X$ 的底空间离散，$X$ 只有有限多个点，设为 $x_1,\dots,x_n$，且每个 $\{x_i\}$ 都是 $X$ 的开子集。可以验证 $X=\Spec\bigl(\prod_{i=1}^n\cO_{X,x_i}\bigr)$，且 $\prod_{i=1}^n\cO_{X,x_i}$ 是阿廷环。

(iii)$\Rightarrow$(i) 由 [Atiyah-Macdonald] 的 2.8.1 与 2.8.3 推出。
\end{proof}

\begin{proposition}\label{it:2.5.12}
设 $f\colon X\to Y$ 是诺特概形的态射。则下列条件等价：
\begin{enumerate}[label=(\roman*)]
\item $f$ 是有限态射；
\item $f$ 是固有且仿射的；
\item $f$ 是固有且拟有限的。
\end{enumerate}
\end{proposition}

\begin{proof}
设 $f$ 是有限态射。则按定义 $f$ 是仿射的，这说明 $f$ 是分离的。对任何 $y\in Y$，投影
\[
f^{-1}(y)=X\times_Y\Spec k(y)\to\Spec k(y)
\]
是有限态射，从而 $X\times_Y\Spec k(y)=\Spec A$，其中 $A$ 是某个作为 $k(y)$-模有限生成的 $k(y)$-代数。但这时 $A$ 是阿廷的，故由 \ref{it:2.5.11}，$f$ 拟有限。下面证明 $f$ 在底拓扑空间上是闭映射。这是关于 $Y$ 的局部问题，故可设 $Y=\Spec B$ 是仿射的。则 $X=\Spec C$，其中 $C$ 是某个作为 $B$-模有限生成的 $B$-代数。设 $\phi\colon B\to C$ 是对应于 $f$ 的同态。对 $C$ 的任何理想 $\mathfrak{c}$，用 [Atiyah-Macdonald] 5.10 可以证明 $f(V(\mathfrak{c}))=V(\phi^{-1}(\mathfrak{c}))$。故 $f$ 是闭态射。由于有限态射的任何基变换都有限，$f$ 是泛闭的。故 $f$ 固有。这就证明了 (i)$\Rightarrow$(ii) 与 (i)$\Rightarrow$(iii)。

设 $f$ 固有。设
\[
X\xrightarrow{f'}Y'\xrightarrow{g}Y
\]
是 $f$ 的 Stein 分解。若 $f$ 是仿射的，则 $f'$ 必是同构。于是 $f$ 可与 $g$ 等同，从而是有限态射。这证明了 (ii)$\Rightarrow$(i)。

若 $f$ 拟有限，则对每个 $y'\in Y'$，$f^{-1}(g(y'))$ 是离散的，从而 $f'^{-1}(y')\subset f^{-1}(g(y'))$ 是离散的。但 $f'^{-1}(y')$ 非空且连通。故 $f'^{-1}(y')$ 只含一个点。于是 $f'$ 是单且满的。此外，由 \ref{it:1.3.26} (vi)，$f'$ 是固有的，故它在底拓扑空间上是同胚。结合 $f'^{\sharp}\colon\cO_{Y'}\to f'_*\cO_X$ 是同构这一事实，我们看到 $f'$ 是概形的同构。于是 $f$ 可与 $g$ 等同，从而是有限态射。这证明了 (iii)$\Rightarrow$(i)。
\end{proof}

\begin{proposition}\label{it:2.5.13}
设 $f\colon X\to Y$ 是诺特概形的固有态射，$X'$ 是 $X$ 中那些在 $f^{-1}(f(x))$ 中孤立的点 $x$ 构成的子集。则 $X'$ 是 $X$ 的开子集。设
\[
X\xrightarrow{f'}Y'\xrightarrow{g}Y
\]
是 $f$ 的 Stein 分解。则 $f'|_{X'}\colon X'\to Y'$ 是开浸入。
\end{proposition}

\begin{proof}
设 $x$ 是 $X'$ 中的点。由于 $f'^{-1}(f'(x))\subset f^{-1}(f(x))$，$x$ 在 $f'^{-1}(f'(x))$ 中是孤立的。我们可以在
\[
f'^{-1}(f'(x))\cong X\times_{Y'}\Spec k(f'(x))
\]
中找到仿射开子概形 $U=\Spec A$，使 $U=\{x\}$，其中 $k(f'(x))$ 是 $\cO_{Y',f'(x)}$ 的剩余域。环 $A$ 必定是阿廷 $k(f'(x))$-代数。特别地，$U\to\Spec k(f'(x))$ 是固有的。由 \ref{it:1.3.26} (vi)，这说明 $U\hookrightarrow f'^{-1}(f'(x))$ 是固有的。故集合 $\{x\}$ 在 $f'^{-1}(f'(x))$ 中既开又闭。但 $f'^{-1}(f'(x))$ 连通，故 $f'^{-1}(f'(x))=\{x\}$。注意由 \ref{it:1.3.26} (vi)，$f'$ 是固有的。故对 $x$ 在 $X$ 中的任何开邻域 $W$，$f'(X-W)$ 是闭的。由于 $f'^{-1}(f'(x))=\{x\}$，$f'(x)$ 落在 $Y'-f'(X-W)$ 中。我们有
\[
f'^{-1}(Y'-f'(X-W))\subset W.
\]
这说明：对 $x$ 在 $X$ 中的任何开邻域 $W$，存在 $f'(x)$ 在 $Y'$ 中的开邻域 $W'=Y'-f'(X-W)$，使 $f'^{-1}(W')\subset W$。这推出
\[
(f'_*\cO_X)_{f'(x)}\cong\cO_{X,x}.
\]
由于 $\cO_{Y'}\cong f'_*\cO_X$，我们有
\[
\cO_{Y',f'(x)}\cong\cO_{X,x}.
\]
由 \ref{it:1.3.13} (iii)，存在 $x$ 在 $X$ 中的邻域 $W_x$，使 $f'|_{W_x}\colon W_x\to Y'$ 是开浸入。对任何 $P\in W_x$，由于 $g$ 是有限态射，$f'(P)$ 在 $g^{-1}(f(P))$ 中是孤立的。故可找到 $f'(P)$ 在 $Y'$ 中的开邻域 $W'_{f'(P)}$，使 $W'_{f'(P)}\cap g^{-1}(f(P))=\{f'(P)\}$。于是我们有
\[
W_x\cap f'^{-1}(W'_{f'(P)})\cap f^{-1}(f(P))=\{P\}.
\]
但 $W_x\cap f'^{-1}(W'_{f'(P)})$ 是 $P$ 的邻域。故 $P$ 在 $f^{-1}(f(P))$ 中是孤立的。于是 $W_x\subset X'$。这说明 $X'$ 在 $X$ 中是开的。此外，存在 $X'$ 的开覆盖 $\{W_x\}_{x\in X'}$，使每个 $f'|_{W_x}\colon W_x\to Y'$ 是开浸入。由于对任何 $x\in X'$，$f'^{-1}(f'(x))=\{x\}$，$f'|_{X'}$ 是单的。故 $f'|_{X'}$ 是开浸入。
\end{proof}

\begin{tnote}
原文此处作「The ring $A$ must an artinian $k(f'(x))$-algebra.」，脱漏 be 动词；译文按「must be」理解。中文没有对应的形态变化，故正文不作改动，仅记于此。[依据：原书 p.207 命题 2.5.13 的证明]
\end{tnote}

\begin{theorem}\label{it:2.5.14}
（\emph{Zariski 主定理}）设 $f\colon X\to Y$ 是诺特概形的拟射影拟有限态射。则我们可以把 $f$ 分解为
\[
X\xrightarrow{j}\ol{X}\xrightarrow{\bar f}Y,
\]
其中 $j$ 是开浸入，$\bar f$ 是有限态射。
\end{theorem}

\begin{proof}
由拟射影态射的定义与 \ref{it:1.4.17}，我们可以把 $f$ 分解为
\[
X\xrightarrow{j'}\ol{X}'\xrightarrow{f'}Y,
\]
其中 $j'$ 是开浸入，$f'$ 是射影态射。考虑 $f'$ 的 Stein 分解
\[
\ol{X}'\xrightarrow{f''}\ol{X}\to Y.
\]
由于 $f$ 拟有限，每个 $x\in X$ 在 $f'^{-1}(f'(x))$ 中是孤立的。由 \ref{it:2.5.13}，$f''|_X\colon X\to\ol{X}$ 是开浸入。$f$ 的分解
\[
X\xrightarrow{f''|_X}\ol{X}\to Y
\]
具有所需性质。
\end{proof}

\begin{corollary}\label{it:2.5.15}
设 $Y$ 是诺特整概形，且对每个 $y\in Y$，整环 $\cO_{Y,y}$ 是整闭的（具有此性质的概形称为\emph{正规}\idx{zhenggui}{正规}的），并设 $X$ 是诺特整概形。则每个分离的拟有限双有理态射 $f\colon X\to Y$ 都是开浸入。
\end{corollary}

\begin{proof}
我们首先证明 $f$ 局部是开浸入。这是局部问题。可设 $X=\Spec B$ 与 $Y=\Spec A$ 都是仿射的，其中 $B$ 是有限生成 $A$-代数。设 $b_1,\dots,b_m$ 是 $B$ 的一族有限生成元。考虑 $A$-代数同态
\[
A[x_1,\dots,x_m]\to B,\qquad x_i\mapsto b_i.
\]

它是满的，从而诱导闭浸入
\[
X=\Spec B\to\Spec A[x_1,\dots,x_m].
\]
由于我们有同构
\[
A[x_0,\dots,x_m]_{(x_0)}\to A[x_1,\dots,x_m],\qquad \frac{x_i}{x_0}\mapsto x_i,
\]
我们得到开浸入
\[
\Spec A[x_1,\dots,x_m]\cong D_+(x_0)\to\bP^m_A.
\]
于是我们得到浸入 $X\to\bP^m_A$。这说明 $X\to Y$ 是拟射影的。由 Zariski 主定理 \ref{it:2.5.14}，$f$ 可分解为
\[
X\xrightarrow{j}\ol{X}\xrightarrow{\bar f}Y,
\]
其中 $j$ 是开浸入，$\bar f$ 是有限态射。用 $j$ 的概形论像替换 $\ol{X}$，可设 $\ol{X}$ 是整的且 $X$ 在 $\ol{X}$ 中稠密。由于 $\bar f$ 有限且双有理、$Y$ 正规，$\bar f$ 必是同构。（参见 [Atiyah-Macdonald] 第 5 章。）故 $f$ 可与 $j$ 等同，从而是开浸入。

为证明 $f$ 是开浸入，我们只需证明 $f$ 作为底拓扑空间上的映射是单的。设 $f(x_1)=f(x_2)=y$。设 $\xi$ 是 $X$ 的一般点，$\eta$ 是 $Y$ 的一般点。由于 $f$ 双有理，我们有 $f(\xi)=\eta$，且 $f^{\sharp}_{\xi}\colon\cO_{Y,\eta}\to\cO_{X,\xi}$ 是同构。注意由于 $x_i\in\ol{\{\xi\}}$，对每个 $i\in\{1,2\}$，$x_i$ 的每个邻域 $U$ 也是 $\xi$ 的邻域，从而我们有同态 $\cO_X(U)\to\cO_{X,\xi}$。取有向极限，我们得到同态
\[
\cO_{X,x_i}=\dlim_{x_i\in U}\cO_X(U)\to\cO_{X,\xi}.
\]
由于 $X$ 是整的，这个同态把 $\cO_{X,\xi}$ 与 $\cO_{X,x_i}$ 的分式域等同，从而是单射。类似地，我们有单射 $\cO_{Y,y}\to\cO_{Y,\eta}$。下图交换：
\[
\begin{tikzcd}
\cO_{Y,y}\arrow[r]\arrow[d,"f^{\sharp}_{x_i}"'] & \cO_{Y,\eta}\arrow[d,"f^{\sharp}_{\xi}"]\\
\cO_{X,x_i}\arrow[r] & \cO_{X,\xi}
\end{tikzcd}
\]
由于 $f$ 局部是开浸入，竖直箭头是同构。由于水平箭头是单射，必有 $\cO_{X,x_1}=\cO_{X,x_2}$。这里我们通过单射 $\cO_{X,x_i}\to\cO_{X,\xi}$ 把 $\cO_{X,x_i}$（$i=1,2$）看作 $\cO_{X,\xi}$ 的子环。设 $U_i$（$i=1,2$）是 $x_i$ 的仿射开邻域，使 $f(U_i)$ 都含在 $y$ 在 $Y$ 中的同一个仿射开邻域 $V$ 中。由 \ref{it:1.3.25} (i)，

$U_1\cap U_2$ 是仿射的，且作为 $\cO_Y(V)$-代数，$\cO_X(U_1\cap U_2)$ 由限制映射 $\cO_X(U_1)\to\cO_X(U_1\cap U_2)$ 与 $\cO_X(U_2)\to\cO_X(U_1\cap U_2)$ 的像生成。我们有交换图
\[
\begin{tikzcd}
& \cO_X(U_1)\arrow[dl]\arrow[dr] &\\
\cO_X(U_1\cap U_2)\arrow[rr] & & \cO_{X,\xi}\\
& \cO_X(U_2)\arrow[ul]\arrow[ur] &
\end{tikzcd}
\]
其中所有箭头都是单射。故我们可以把 $\cO_X(U_1)$、$\cO_X(U_2)$ 与 $\cO_X(U_1\cap U_2)$ 看作 $\cO_{X,\xi}$ 的子代数。此外，$\cO_X(U_1\cap U_2)$ 是 $\cO_{X,\xi}$ 中由 $\cO_X(U_1)$ 与 $\cO_X(U_2)$ 生成的子代数。我们有交换图
\[
\begin{tikzcd}
& \cO_X(U_1)\arrow[dl]\arrow[dr] &\\
\cO_{X,x_1}=\cO_{X,x_2}\arrow[rr] & & \cO_{X,\xi}\\
& \cO_X(U_2)\arrow[ul]\arrow[ur] &
\end{tikzcd}
\]
其中所有箭头都是单射。故若把 $\cO_{X,x_1}=\cO_{X,x_2}$ 看作 $\cO_{X,\xi}$ 的子代数，则它包含 $\cO_{X,\xi}$ 中由 $\cO_X(U_1)$ 与 $\cO_X(U_2)$ 生成的子代数，即
\[
\cO_X(U_1\cap U_2)\subset\cO_{X,x_1}=\cO_{X,x_2}.
\]
设 $\fm_x$ 是 $\cO_{X,x_1}=\cO_{X,x_2}$ 的极大理想。则 $\fm_x\cap\cO_X(U_1\cap U_2)$ 是 $\cO_X(U_1\cap U_2)$ 的素理想。由于 $U_1\cap U_2$ 仿射，它对应于 $U_1\cap U_2$ 中的点 $x$。对每个 $i$，$\cO_X(U_i)$ 的素理想 $\fm_x\cap\cO_X(U_i)$ 对应于仿射概形 $U_i$ 中的点 $x_i$。在包含 $U_1\cap U_2\hookrightarrow U_i$ 下，$x$ 映到 $x_i$。故 $x_1=x_2=x$，从而 $f$ 是单射。
\end{proof}

\begin{tnote}
原文把 $D_+(x_0)$ 上的坐标 $x_i/x_0$ 简记为 $x_i$，故同构右端的 $\Spec A[x_1,\dots,x_m]$ 实为
\[
\Spec A[x_1/x_0,\dots,x_m/x_0].
\]
而 $A[x_0,\dots,x_m]_{(x_0)}$ 是分次环在齐次素理想 $(x_0)$ 处的齐次局部化，其定义见原书 p.28。[依据：原书 p.209 推论 2.5.15 的证明]
\end{tnote}
